\documentclass[10pt]{article}
\usepackage[a4paper,margin=13mm]{geometry}
\usepackage{graphicx}
\usepackage{xcolor}
\usepackage{helvet}
\renewcommand{\familydefault}{\sfdefault}
\pagestyle{empty}
\newcommand{\panel}[3]{%
  \begin{minipage}[t]{#1}
  \textbf{\Large #2}\vspace{1mm}\\
  \includegraphics[width=\linewidth]{#3}
  \end{minipage}
}
\begin{document}
\noindent{\Large\textbf{pyccc: CellChat-compatible CCC at AnnData scale}}\\[2mm]
\noindent\panel{0.47\textwidth}{a}{subfigures/panel_a_workflow.pdf}\hfill
\panel{0.47\textwidth}{b}{subfigures/panel_b_parity.pdf}\\[5mm]
\noindent\panel{0.47\textwidth}{c}{subfigures/panel_c_runtime.pdf}\hfill
\panel{0.47\textwidth}{d}{subfigures/panel_d_summary_tradeoff.pdf}\\[4mm]
{\small
\textbf{Figure 1.} pyccc keeps CellChat-like CCC in an AnnData-native sparse workflow,
matches CellChat R numerics on the official tutorial benchmark, accelerates the
real 1M-cell workflow by avoiding CellChat's dense cell-level path, and exposes
outlier-robust summary options with explicit similarity tradeoffs.
}
\end{document}
